\chapter{Urine Protein-to-Creatinine Ratio}

\section*{What Is This Test?}
The urine protein-to-creatinine ratio uses a single urine sample to estimate how much protein is lost in the urine over a full day. It roughly adjusts for how concentrated or dilute the urine is. It offers a convenient alternative to a timed urine collection for quantifying proteinuria.

\section*{Alternative Names}
\begin{itemize}
  \item UPCR
  \item Protein/Creatinine Ratio
  \item Random Urine Protein-to-Creatinine Ratio
\end{itemize}
\section*{Why This Test Is Ordered}
It is used to evaluate and monitor proteinuria in chronic kidney disease, glomerular disease, diabetes, hypertension, pregnancy-related hypertensive disorders, and systemic diseases that affect the kidneys. The albumin-to-creatinine ratio remains preferred when the focus is specifically on albuminuria and on cardiovascular or kidney risk.

\section*{How This Test Is Performed}
A clean-catch spot urine specimen is collected, preferably under consistent conditions. The laboratory measures total protein and creatinine and reports their ratio, commonly in mg/g or mg/mmol.

\section*{Preparation, Risks, and Limitations}
No fasting is generally required. The test is noninvasive. Fever, strenuous exercise, urinary infection, blood in the urine, heart failure, menstruation, and highly concentrated or dilute urine can cause transient or misleading elevations. An unexpected abnormal result is usually repeated.

\section*{Understanding Results}
Higher ratios indicate greater protein loss and may suggest glomerular kidney injury. Persistent elevation requires confirmation and evaluation with eGFR, urine sediment, blood pressure, diabetes assessment, and the clinical context. Thresholds vary by laboratory and guideline.

\section*{References}
Kidney Disease: Improving Global Outcomes guidance on proteinuria and albuminuria assessment.

\clearpage
\section*{Understanding Results and Follow-up}
The laboratory report should identify the specimen, method, units, reference interval or decision limit, and whether the result is positive, negative, abnormal, or inconclusive. Reference intervals vary with age, sex, pregnancy, specimen type, assay, and laboratory; a result outside a range is not automatically a diagnosis, and a result within a range does not always exclude disease. Discuss unexpected results with the ordering clinician, who can decide whether repeat or confirmatory testing, treatment, or referral is appropriate. Do not change medicines or delay urgent care based on an isolated result.


\section*{Regulatory Status and Authorization Date}
This chapter describes a test category, not a specific commercial product. FDA, CDSCO, EU IVDR, and MHRA approval or registration dates therefore cannot be assigned without the assay manufacturer, product name, instrument, and intended use. For laboratory-developed tests, an individual agency approval date may not exist. See the Preface for the interpretation of regulatory status.

\clearpage
